\documentclass[uplatex,a4paper,11pt]{jsarticle}
\usepackage[margin=22mm]{geometry}
\usepackage[dvipdfmx]{hyperref,xcolor}
\usepackage{fvextra,enumitem,longtable,booktabs,amssymb}
\usepackage{listings,jvlisting}
\hypersetup{colorlinks=true,linkcolor=blue!55!black,urlcolor=blue!65!black,
  pdftitle={P2 TeX・LilyPond・移調サービス構築手順}}
\DefineVerbatimEnvironment{command}{Verbatim}{fontsize=\small,frame=single,
  rulecolor=\color{gray!55},breaklines=true,breakanywhere=true}
\setlist{nosep}
\lstset{language=Python,basicstyle=\ttfamily\scriptsize,numbers=left,
  numberstyle=\tiny,frame=single,breaklines=true,columns=fullflexible,
  keepspaces=true,showstringspaces=false,tabsize=4,captionpos=b}

\title{P2 TeX・LilyPond・移調\\サービス構築手順}
\author{}
\date{2026年9月6日}

\begin{document}
\maketitle
\begin{abstract}
本書は、同じP2上でTeX PDF生成、LilyPond PDF生成、LilyPond移調テキスト生成の
3 APIを別プロセスとして構築し、
P1からの内部トークン付き要求だけを受け付ける手順である。P2は利用者のパスワード、
ユーザーDB、JWT署名鍵を保持せず、P0やクライアントから直接接続させない。
\end{abstract}
\tableofcontents
\newpage

\section{最初に理解する完成構成}

\begin{command}
P1 TeX認証API       -- token T --> P2 TeX API       TCP 8000
P1 LilyPond認証API  -- token L --> P2 LilyPond PDF API       TCP 8001
                         `-------> P2 LilyPond移調API         TCP 8002
\end{command}

内部トークンTとLは別の値である。LilyPondの二つのAPIは同じ内部トークンLを使う。
P2のUFWとVLAN ACLはP1固定IPから3ポートへの通信だけ
を許可する。P2の\texttt{/health}は認証不要のため、ファイアウォール境界の外へ公開しない。

\section{前提値とディレクトリ}

\begin{longtable}{@{}p{48mm}p{50mm}p{48mm}@{}}
\toprule
項目 & 例 & 用途\\
\midrule
P1固定IP & \texttt{192.168.30.10} & 接続を許可する唯一のホスト\\
P2固定IP & \texttt{192.168.40.10} & コンパイルVLAN\\
管理ネットワーク & \texttt{192.168.99.0/24} & SSH管理\\
TeX API & \path{/opt/texnote/P2}, TCP 8000 & 5種類のTeXエンジン\\
LilyPond PDF API & \path{/opt/lilypondnote/P2}, TCP 8001 & LilyPond PDF生成\\
LilyPond移調API & \path{/opt/lilypond-transpose/P2}, TCP 8002 & 移調後テキスト生成\\
\bottomrule
\end{longtable}

移調サービスの配置先は次で固定する。\texttt{engine}には
LilyPond\_transposeリポジトリのPython実装だけを置き、Web APIラッパー、venv、エンジンを
一つの専用ルート以下へ分離する。

\begin{command}
/opt/lilypond-transpose/P2/
|-- transpose_app.py          HTTP APIラッパー
|-- requirements-transpose.txt
|-- .venv/                    移調API専用Python環境
`-- engine/                   LilyPond_transposeの実処理
    |-- main.py               CLIエントリポイント
    |-- music_theory.py
    |-- pitch.py
    |-- relative.py
    |-- tokenizer.py
    |-- tokens.py
    |-- transpose.py
    `-- writer.py
\end{command}

systemdは\texttt{WorkingDirectory=/opt/lilypond-transpose/P2}から
\texttt{transpose\_app:app}を起動し、ラッパーは
\path{/opt/lilypond-transpose/P2/engine/main.py}を固定コマンドとして呼ぶ。利用者から実行
ファイルのパスを指定させない。

\section{OSとコンパイラを導入する}

64-bit Raspberry Pi OS LiteまたはDebianを導入し、固定IPを設定する。

\begin{command}
sudo apt update
sudo apt full-upgrade
sudo apt install python3-venv python3-pip ufw \
  texlive-latex-base texlive-latex-extra texlive-luatex \
  texlive-xetex texlive-lang-japanese latexmk lilypond
\end{command}

必要な実行ファイルを確認する。

\begin{command}
for cmd in lualatex xelatex pdflatex uplatex platex dvipdfmx \
  extractbb lilypond; do command -v "$cmd" || echo "MISSING: $cmd"; done
lilypond --version
\end{command}

文書が追加フォントやTeXパッケージを使用する場合は、実際の入力で単体試験してから不足
パッケージを追加する。サーバープロセスから\texttt{tlmgr}を実行させない。

\section{専用ユーザーと配置先を作る}

\begin{command}
sudo useradd --system --home /opt/texnote/P2 \
  --shell /usr/sbin/nologin texnote-p2
sudo useradd --system --home /opt/lilypondnote/P2 \
  --shell /usr/sbin/nologin lilypondnote-p2
sudo useradd --system --home /opt/lilypond-transpose/P2 \
  --shell /usr/sbin/nologin lilypond-transpose-p2
sudo install -d -o texnote-p2 -g texnote-p2 /opt/texnote/P2
sudo install -d -o lilypondnote-p2 -g lilypondnote-p2 /opt/lilypondnote/P2
sudo install -d -o lilypond-transpose-p2 -g lilypond-transpose-p2 \
  /opt/lilypond-transpose/P2
\end{command}

3サービスを別OSユーザーで動かす。PDF生成サービスの一時作業はPythonの一時ディレクトリ内
で行われ、完了後に削除される。

\section{TeX APIを配置する}

管理用MacのTeXNoteリポジトリから転送する。

\begin{command}
# 管理用Macで実行
scp Server/P2/progs/app.py Server/P2/progs/requirements.txt \
  mat@192.168.40.10:/tmp/
\end{command}

P2で配置し、仮想環境を作る。

\begin{command}
sudo install -o texnote-p2 -g texnote-p2 -m 640 \
  /tmp/app.py /tmp/requirements.txt /opt/texnote/P2/
sudo -u texnote-p2 python3 -m venv /opt/texnote/P2/.venv
sudo -u texnote-p2 /opt/texnote/P2/.venv/bin/pip install \
  -r /opt/texnote/P2/requirements.txt
sudo -u texnote-p2 /opt/texnote/P2/.venv/bin/python \
  -m py_compile /opt/texnote/P2/app.py
\end{command}

\subsection{TeX側の必要ソース}

ここで配置した二つのファイルが、P2のTeX PDF生成サービスに必要なPython側の全ファイルで
ある。\texttt{requirements.txt}は\path{/opt/texnote/P2/.venv}だけに導入するPython
パッケージを指定している。

\lstinputlisting[caption={TeX PDF生成 app.py}]{/Users/mat/Documents/Git/TeXNote/Server/P2/progs/app.py}
\lstinputlisting[language={},caption={TeX API専用 requirements.txt}]{/Users/mat/Documents/Git/TeXNote/Server/P2/progs/requirements.txt}

\section{LilyPond APIを配置する}

管理用MacのLilyPondNoteリポジトリから転送する。

\begin{command}
# 管理用Macで実行
scp Server/P2/app.py Server/P2/requirements.txt \
  mat@192.168.40.10:/tmp/
\end{command}

同名ファイルをTeX用とは別ディレクトリへ配置する。

\begin{command}
sudo install -o lilypondnote-p2 -g lilypondnote-p2 -m 640 \
  /tmp/app.py /tmp/requirements.txt /opt/lilypondnote/P2/
sudo -u lilypondnote-p2 python3 -m venv /opt/lilypondnote/P2/.venv
sudo -u lilypondnote-p2 /opt/lilypondnote/P2/.venv/bin/pip install \
  -r /opt/lilypondnote/P2/requirements.txt
sudo -u lilypondnote-p2 /opt/lilypondnote/P2/.venv/bin/python \
  -m py_compile /opt/lilypondnote/P2/app.py
\end{command}

\subsection{LilyPond PDF生成側の必要ソース}

ここで配置した\texttt{app.py}はLilyPondを実行してPDFを返すAPIであり、後述する移調API
とは別プロセスである。\texttt{requirements.txt}は
\path{/opt/lilypondnote/P2/.venv}だけに導入する。

\lstinputlisting[caption={LilyPond PDF生成 app.py}]{/Users/mat/Documents/Git/LilyPondNote/Server/P2/app.py}
\lstinputlisting[language={},caption={LilyPond PDF生成API専用 requirements.txt}]{/Users/mat/Documents/Git/LilyPondNote/Server/P2/requirements.txt}

\section{LilyPond移調テキストAPIを配置する}

\subsection{CLIプログラムの転送}

LilyPond\_transposeリポジトリの現行プログラムはWeb APIではない。
\path{python/main.py}を\texttt{SRC DST}引数で実行すると、標準入力から楽譜テキストを読み、移調後テキストを
標準出力へ返すCLIである。サンプル楽譜、生成済みPDF、シェルスクリプトはサーバー実行に
不要なので転送しない。

\begin{command}
# 管理用MacでLilyPond_transposeリポジトリ直下から実行
ssh mat@192.168.40.10 'mkdir -p /tmp/lilypond-transpose'
scp python/main.py python/music_theory.py python/pitch.py \
  python/relative.py python/tokenizer.py python/tokens.py \
  python/transpose.py python/writer.py \
  mat@192.168.40.10:/tmp/lilypond-transpose/
\end{command}

P2では\path{/opt/lilypond-transpose/P2/engine}へ配置する。このディレクトリを移調
アルゴリズムの正式な運転場所とする。\path{/usr/local/bin}や管理ユーザーのホームへ置かない。

\begin{command}
sudo install -d -o lilypond-transpose-p2 -g lilypond-transpose-p2 \
  /opt/lilypond-transpose/P2/engine
sudo install -o lilypond-transpose-p2 -g lilypond-transpose-p2 -m 640 \
  /tmp/lilypond-transpose/*.py /opt/lilypond-transpose/P2/engine/
sudo -u lilypond-transpose-p2 python3 -m py_compile \
  /opt/lilypond-transpose/P2/engine/*.py
\end{command}

代表的な入力でCLI単体試験を行う。

\begin{command}
sudo -u lilypond-transpose-p2 sh -c \
  'printf "%s\n" "\\relative c'\'' { c4 d e f }" | \
   python3 /opt/lilypond-transpose/P2/engine/main.py c a'
\end{command}

\subsection{移調CLI側の必要ファイル}

次の8ファイルは、直前の手順で\path{/opt/lilypond-transpose/P2/engine}へ配置する
移調エンジン一式である。\texttt{main.py}だけでは動作しないため、8ファイルを同じ
ディレクトリへ配置する。サンプル楽譜、PDF、生成用シェルスクリプトはP2の運転には
使用しないので掲載しない。8ファイルの内容は、本文の構築手順を中断しないよう
付録「Python移調プログラムのソース」にまとめて掲載する。

\subsection{APIラッパーの必須仕様}

CLIをHTTP経由で安全に呼び出すため、LilyPondNoteリポジトリの
\path{Server/P2/transpose_app.py}を使用する。API契約は次のとおりである。

\begin{command}
POST /v1/transpose
Authorization: Bearer <LILYPONDNOTE_API_TOKEN>
Content-Type: application/json

request:
{
  "sourcePitch": "g",
  "destinationPitch": "d'",
  "source": "LilyPond source text"
}

response:
{
  "transposedSource": "transposed LilyPond source text"
}
\end{command}

このAPIは\textbf{LilyPondコマンドを実行せず、PDFを生成・返却しない}。内部では固定された
Python実行ファイルと\texttt{main.py}を引数リストで起動し、\texttt{source}を標準入力へ
渡す。シェル、リダイレクト文字列、ユーザー指定コマンドを実行してはならない。

\texttt{sourcePitch}と\texttt{destinationPitch}には、CLIが対応するLilyPond音名と、任意の
オクターブ指示を指定できる。ASCIIアポストロフィー\texttt{'}は1個につき1オクターブ上、
カンマ\texttt{,}は1個につき1オクターブ下を意味する。例えば\texttt{g}から\texttt{d}は
5半音下、\texttt{g}から\texttt{d'}は7半音上、\texttt{g}から\texttt{d,}は17半音下である。
JSON文字列として送信し、P2はシェルを介さず引数配列へ渡すため、APIデータ内の
アポストロフィーにバックスラッシュを付けない。\texttt{d\textbackslash'}は無効である。

アポストロフィーとカンマは同じ指示内で混在させない。\texttt{c''}や\texttt{c,,}は有効、
\texttt{c',}は無効である。P2 APIが受理するのはASCIIアポストロフィーだけであり、スマート
引用符はLilyPondNoteアプリ側でASCIIへ正規化してから送信する。APIを直接呼ぶクライアントは
最初からASCIIアポストロフィーを使用する。

ラッパーには次の制限を実装する。

\begin{itemize}
  \item Pydanticで未知フィールドを拒否し、入力テキストを2,000,000文字以下にする。
  \item 移調元・移調先を「CLIが対応する音名＋同種のアポストロフィーまたはカンマ」で検証する。
  \item subprocessへタイムアウト、CPU、メモリ、出力サイズの上限を設定する。
  \item 終了コードが0でない場合はPDF生成APIと区別できるエラーとして返す。
  \item ログ、例外、journalへ入力楽譜全文や内部トークンを出力しない。
  \item 応答はUTF-8の\texttt{transposedSource}だけとし、\texttt{pdfBase64}を含めない。
\end{itemize}

ラッパー、専用requirements、完成済みsystemdファイル、env雛形はLilyPondNoteの
\path{Server/P2}で一緒に管理する。管理用Macから転送し、専用venvへrequirementsに記載した
FastAPIとUvicornだけを導入する。

\begin{command}
# 管理用MacでLilyPondNoteリポジトリ直下から実行
scp Server/P2/transpose_app.py Server/P2/requirements-transpose.txt \
  Server/P2/lilypond-transpose-p2.service \
  Server/P2/lilypond-transpose-p2.env \
  mat@192.168.40.10:/tmp/

# P2で実行
sudo install -o lilypond-transpose-p2 -g lilypond-transpose-p2 -m 640 \
  /tmp/transpose_app.py /tmp/requirements-transpose.txt \
  /opt/lilypond-transpose/P2/
sudo -u lilypond-transpose-p2 python3 -m venv \
  /opt/lilypond-transpose/P2/.venv
sudo -u lilypond-transpose-p2 \
  /opt/lilypond-transpose/P2/.venv/bin/pip install \
  --requirement /opt/lilypond-transpose/P2/requirements-transpose.txt
sudo -u lilypond-transpose-p2 \
  /opt/lilypond-transpose/P2/.venv/bin/python -m py_compile \
  /opt/lilypond-transpose/P2/transpose_app.py \
  /opt/lilypond-transpose/P2/engine/*.py
\end{command}

\subsection{移調API側の必要ソース}

\texttt{transpose\_app.py}はHTTP要求を検証し、直前に掲載した
\texttt{engine/main.py}を固定コマンドとして呼び出すラッパーである。この専用venvに必要な
外部Pythonパッケージは、次のrequirementsにあるFastAPIとUvicornだけである。

\lstinputlisting[caption={移調テキストAPI transpose\_app.py}]{/Users/mat/Documents/Git/LilyPondNote/Server/P2/transpose_app.py}
\lstinputlisting[language={},caption={移調API専用 requirements-transpose.txt}]{/Users/mat/Documents/Git/LilyPondNote/Server/P2/requirements-transpose.txt}

\section{内部トークンを設定する}

トークンはP1で生成した値を安全な経路でP2へ設定する。TeX用とLilyPond用を取り違えない。
Mac側には次の3個のenv雛形がある。

\begin{itemize}
  \item TeX用：TeXNoteリポジトリの\path{Server/P2/p2.env}
  \item LilyPond PDF用：LilyPondNoteリポジトリの
        \path{Server/P2/lilypondnote-p2.env}
  \item LilyPond移調用：LilyPondNoteリポジトリの
        \path{Server/P2/lilypond-transpose-p2.env}
\end{itemize}

TeX用雛形で\texttt{TEXNOTE\_API\_TOKEN}の右辺にある山括弧内の仮の値は、P1の
\path{/etc/texnote/p1.env}にある\texttt{TEXNOTE\_API\_TOKEN}と同じ値へ置き換える。
LilyPond側二つの\texttt{<...>}は、P1の\path{/etc/lilypondnote/p1.env}にある
\texttt{LILYPONDNOTE\_API\_TOKEN}と同じLトークンへ置き換える。TとLは別の値である。

\begin{command}
# /etc/texnote/p2.env
TEXNOTE_API_TOKEN=<P1 TeX側と同じ内部トークンT>

# /etc/lilypondnote/p2.env
LILYPONDNOTE_API_TOKEN=<P1 LilyPond側と同じ内部トークンL>
LILYPOND_EXECUTABLE=/usr/bin/lilypond

# /etc/lilypond-transpose/p2.env
LILYPONDNOTE_API_TOKEN=<P1 LilyPond側と同じ内部トークンL>
LILYPOND_TRANSPOSE_MAIN=/opt/lilypond-transpose/P2/engine/main.py
\end{command}

次の完成形の雛形をMac側で編集する。秘密値の\texttt{<...>}だけを上記の対応するトークンへ
置き換え、変数名、LilyPond実行ファイル、移調エンジンの配置先は変更しない。

\lstinputlisting[language={},caption={TeX PDF生成用 p2.env雛形}]{/Users/mat/Documents/Git/TeXNote/Server/P2/p2.env}
\lstinputlisting[language={},caption={LilyPond PDF生成用 p2.env雛形}]{/Users/mat/Documents/Git/LilyPondNote/Server/P2/lilypondnote-p2.env}
\lstinputlisting[language={},caption={LilyPond移調用 p2.env雛形}]{/Users/mat/Documents/Git/LilyPondNote/Server/P2/lilypond-transpose-p2.env}

\begin{command}
# 管理用MacのTeXNoteリポジトリ直下で実行
scp Server/P2/p2.env mat@192.168.40.10:/tmp/texnote-p2.env

# 管理用MacのLilyPondNoteリポジトリ直下で実行
scp Server/P2/lilypondnote-p2.env \
  mat@192.168.40.10:/tmp/lilypondnote-p2.env
scp Server/P2/lilypond-transpose-p2.env \
  mat@192.168.40.10:/tmp/lilypond-transpose-p2.env

# P2で実行：存在しない/etc配下を先に作る
sudo install -d -o root -g texnote-p2 -m 750 /etc/texnote
sudo install -o root -g texnote-p2 -m 640 \
  /tmp/texnote-p2.env /etc/texnote/p2.env
sudo install -d -o root -g lilypondnote-p2 -m 750 /etc/lilypondnote
sudo install -o root -g lilypondnote-p2 -m 640 \
  /tmp/lilypondnote-p2.env /etc/lilypondnote/p2.env
sudo install -d -o root -g lilypond-transpose-p2 -m 750 \
  /etc/lilypond-transpose
sudo install -o root -g lilypond-transpose-p2 -m 640 \
  /tmp/lilypond-transpose-p2.env /etc/lilypond-transpose/p2.env
\end{command}

配置後に、秘密値そのものを画面へ出さず、ファイルの所有者・権限と変数名だけを確認する。

\begin{command}
sudo stat -c '%U:%G %a %n' /etc/texnote/p2.env \
  /etc/lilypondnote/p2.env /etc/lilypond-transpose/p2.env
sudo sed -n 's/=.*$/=<設定済み>/p' /etc/texnote/p2.env \
  /etc/lilypondnote/p2.env /etc/lilypond-transpose/p2.env
\end{command}

実機では次の結果になった。所有者が\texttt{root}、グループが各サービス専用グループ、
権限が\texttt{640}であり、必要な変数名もすべて存在している。この結果は正常である。

\begin{command}
root:texnote-p2 640 /etc/texnote/p2.env
root:lilypondnote-p2 640 /etc/lilypondnote/p2.env
root:lilypond-transpose-p2 640 /etc/lilypond-transpose/p2.env

TEXNOTE_API_TOKEN=<設定済み>
LILYPONDNOTE_API_TOKEN=<設定済み>
LILYPOND_EXECUTABLE=<設定済み>
LILYPONDNOTE_API_TOKEN=<設定済み>
LILYPOND_TRANSPOSE_MAIN=<設定済み>
\end{command}

ただし、上の\texttt{sed}は変数名の存在だけを確認するため、値が空、または山括弧付きの
雛形のままでも\texttt{<設定済み>}と表示する。値を画面へ出さずに、空欄と未置換の雛形が
ないことを次で確認する。

\begin{command}
sudo awk -F= '
  /^[[:space:]]*#/ || /^[[:space:]]*$/ { next }
  NF < 2 || $2 == "" || $2 ~ /^<.*>$/ {
    print FILENAME ": 未設定または雛形のまま: " $1
    error=1
  }
  END {
    if (!error) print "すべての値が設定済みです"
    exit error
  }
' /etc/texnote/p2.env \
  /etc/lilypondnote/p2.env \
  /etc/lilypond-transpose/p2.env
\end{command}

正常なら\texttt{すべての値が設定済みです}の1行だけが表示される。ファイル名と変数名が
表示された場合は、そのファイルの該当値をMac側の雛形で修正し、再転送・再配置する。

P2にP1のJWT署名鍵を置かない。P2は利用者JWTではなくP1--P2間の固定トークンだけを検証する。

\section{systemdで3サービスを起動する}

\subsection{TeXコンパイルサービス}

TeXNoteリポジトリにある完成済み\path{Server/P2/texnote-p2.service}を使用する。主要行を
P2上で手入力せず、ファイル全体を転送して
\path{/etc/systemd/system/texnote-p2.service}へ配置する。

\begin{command}
# 管理用MacのTeXNoteリポジトリ直下で実行
scp Server/P2/texnote-p2.service mat@192.168.40.10:/tmp/

# P2で実行
sudo install -o root -g root -m 644 /tmp/texnote-p2.service \
  /etc/systemd/system/texnote-p2.service
\end{command}

\lstinputlisting[language={},caption={TeX PDF生成systemdサービス}]{/Users/mat/Documents/Git/TeXNote/Server/P2/texnote-p2.service}

\subsection{LilyPondコンパイルサービス}

リポジトリの完成済み\path{lilypondnote-p2.service}を転送し、固定IPのTCP 8001で
待ち受ける。

\begin{command}
# 管理用MacのLilyPondNoteリポジトリ直下で実行
scp Server/P2/lilypondnote-p2.service mat@192.168.40.10:/tmp/

# P2で実行
sudo install -o root -g root -m 644 /tmp/lilypondnote-p2.service \
  /etc/systemd/system/lilypondnote-p2.service
\end{command}

\lstinputlisting[language={},caption={LilyPond PDF生成systemdサービス}]{/Users/mat/Documents/Git/LilyPondNote/Server/P2/lilypondnote-p2.service}

\subsection{LilyPond移調テキストサービス}

完成済みファイルをMac側から転送する。主要行を手入力せず、リポジトリにある設定全体を
\path{/etc/systemd/system/lilypond-transpose-p2.service}へ配置する。

\begin{command}
# 管理用MacのLilyPondNoteリポジトリ直下で実行
scp Server/P2/lilypond-transpose-p2.service \
  mat@192.168.40.10:/tmp/

# P2で実行
sudo install -o root -g root -m 644 \
  /tmp/lilypond-transpose-p2.service \
  /etc/systemd/system/lilypond-transpose-p2.service
\end{command}

\lstinputlisting[language={},caption={移調テキストsystemdサービス}]{/Users/mat/Documents/Git/LilyPondNote/Server/P2/lilypond-transpose-p2.service}

起動する。

\begin{command}
sudo systemctl daemon-reload
sudo systemctl enable --now texnote-p2.service
sudo systemctl enable --now lilypondnote-p2.service
sudo systemctl enable --now lilypond-transpose-p2.service
systemctl status texnote-p2.service --no-pager -l
systemctl status lilypondnote-p2.service --no-pager -l
systemctl status lilypond-transpose-p2.service --no-pager -l
sudo ss -ltnp | grep -E ':8000|:8001|:8002'
\end{command}

\subsection{envを読み込んで起動できることを確認する}

3サービスの定義とenvを配置した後に一度再起動し、envのファイル名、権限、変数名、値が
サービスから利用できることを確認する。

\begin{command}
sudo systemctl restart texnote-p2.service
sudo systemctl restart lilypondnote-p2.service
sudo systemctl restart lilypond-transpose-p2.service

systemctl status texnote-p2.service --no-pager -l
systemctl status lilypondnote-p2.service --no-pager -l
systemctl status lilypond-transpose-p2.service --no-pager -l
\end{command}

3個とも\texttt{Active: active (running)}であれば、envを読み込んでプロセスを起動できて
いる。失敗したサービスは、\texttt{journalctl -u サービス名 -n 100 --no-pager}で原因を
確認する。なお、ここで確認できるのはenvを読み込んで起動できたことまでである。P1とP2の
トークンが同じ値かは、第11章の正常要求と、トークンなし・取り違え要求の試験で確認する。

\section{P2のUFWを設定する}

SSHとP1からの3サービス経路だけを許可する。

\begin{command}
sudo ufw default deny incoming
sudo ufw default allow outgoing
sudo ufw allow from 192.168.99.0/24 to any port 22 proto tcp \
  comment 'SSH from management network'
sudo ufw allow from 192.168.30.10 to any port 8000 proto tcp \
  comment 'P1 to TeX compiler'
sudo ufw allow from 192.168.30.10 to any port 8001 proto tcp \
  comment 'P1 to LilyPond compiler'
sudo ufw allow from 192.168.30.10 to any port 8002 proto tcp \
  comment 'P1 to LilyPond transpose'
sudo ufw show added
sudo ufw enable
sudo ufw status verbose
\end{command}

別ターミナルからSSHを再接続する。P0、一般LAN、インターネットから8000、8001、8002へ到達
できてはならない。VLAN ACLでも同じ制限を設定し、P2からP0、管理VLAN、家庭内LANへの
新規接続を拒否する。OS更新用DNS、NTP、HTTPSだけを必要に応じて許可する。

\section{P2内とP1から試験する}

P2自身からヘルスチェックする。

\begin{command}
curl -sS http://192.168.40.10:8000/health
curl -sS http://192.168.40.10:8001/health
curl -sS http://192.168.40.10:8002/health
\end{command}

移調APIは\texttt{"status":"ok"}かつ\texttt{"transposeProgram":true}を返すことを確認する。
HTTP 200でも\texttt{unavailable}または\texttt{false}なら、envの
\texttt{LILYPOND\_TRANSPOSE\_MAIN}と\path{/opt/lilypond-transpose/P2/engine/main.py}の
存在・権限が一致していない。

移調APIの結合試験では、オクターブ指示なしだけでなく、\texttt{g}から\texttt{d'}、
\texttt{c}から\texttt{c,}も送信する。前者が7半音上、後者が12半音下の
\texttt{transposedSource}を返すことを確認する。また、\texttt{d',}、スマート引用符を含む
\texttt{d’}、バックスラッシュを含む\texttt{d\textbackslash'}がHTTP 422になることを確認する。

内部トークンなしの処理要求が401になることを確認する。

\begin{command}
curl -i -X POST http://192.168.40.10:8000/v1/typeset \
  -H 'Content-Type: application/json' -d '{}'
curl -i -X POST http://192.168.40.10:8001/v1/typeset \
  -H 'Content-Type: application/json' -d '{}'
curl -i -X POST http://192.168.40.10:8002/v1/transpose \
  -H 'Content-Type: application/json' -d '{}'
\end{command}

次にP1の各認証API経由でTeX PDF、LilyPond PDF、LilyPond移調テキストを生成する。
移調応答に\texttt{transposedSource}があり、\texttt{pdfBase64}がないことを確認する。
P1ログとP2ログを同時に確認する。

\begin{command}
journalctl -u texnote-p2.service -f
journalctl -u lilypondnote-p2.service -f
journalctl -u lilypond-transpose-p2.service -f
\end{command}

内部トークンTをLilyPondへ、トークンLをTeXへ送った場合に401となることも確認する。

\section{資源分離と運用上の注意}

3 APIはプロセス内で同時実行数、CPU時間、仮想メモリ、出力サイズを制限し、
同じP2のCPUとメモリを共有する。メモリ不足や過負荷が起きる場合はsystemdへ
\texttt{MemoryMax}、\texttt{CPUQuota}、\texttt{TasksMax}を追加し、実機のRAM容量と
代表的な文書で値を決める。一律に小さな値を設定すると正常なコンパイルも失敗する。

ログに文書本文、パスワード、トークンを出力しない。更新時は一つずつ再起動し、他の
サービスに影響がないことを確認する。

\begin{command}
sudo systemctl restart texnote-p2.service
sudo systemctl restart lilypondnote-p2.service
sudo systemctl restart lilypond-transpose-p2.service
journalctl -u texnote-p2.service -n 100 --no-pager
journalctl -u lilypondnote-p2.service -n 100 --no-pager
journalctl -u lilypond-transpose-p2.service -n 100 --no-pager
\end{command}

\section{完了チェックリスト}

\begin{itemize}
  \item[$\square$] TeXはTCP 8000、LilyPond PDFは8001、移調テキストは8002で動く。
  \item[$\square$] 3サービスは別OSユーザー、別仮想環境、別環境ファイルを使う。
  \item[$\square$] TeXとLilyPondの内部トークンが異なる。
  \item[$\square$] P2にユーザーDB、パスワード、JWT署名鍵がない。
  \item[$\square$] P1からだけP2の8000、8001、8002へ接続できる。
  \item[$\square$] P0と一般LANからP2へ直接接続できない。
  \item[$\square$] トークンなし・取り違えた要求は401になる。
  \item[$\square$] PDF APIはPDF、移調APIはテキストだけを返す。
  \item[$\square$] 移調APIはLilyPond実行やPDF生成を行わない。
  \item[$\square$] P2再起動後に3サービスとUFWが自動復旧する。
\end{itemize}

\clearpage
\appendix
\section{Python移調プログラムのソース}

ここには、\path{/opt/lilypond-transpose/P2/engine}へ配置するPython移調エンジン8ファイル
だけを掲載する。HTTP APIラッパー、requirements、env、systemdは、使用する手順と対応が
分かるよう本文の各箇所に掲載している。

\subsection{main.py}
\lstinputlisting[caption={移調CLIの入口 main.py}]{/Users/mat/Documents/Git/LilyPond_transpose/python/main.py}

\subsection{music\_theory.py}
\lstinputlisting[caption={音楽理論処理 music\_theory.py}]{/Users/mat/Documents/Git/LilyPond_transpose/python/music_theory.py}

\subsection{pitch.py}
\lstinputlisting[caption={音高処理 pitch.py}]{/Users/mat/Documents/Git/LilyPond_transpose/python/pitch.py}

\subsection{relative.py}
\lstinputlisting[caption={relative記法処理 relative.py}]{/Users/mat/Documents/Git/LilyPond_transpose/python/relative.py}

\subsection{tokenizer.py}
\lstinputlisting[caption={字句解析 tokenizer.py}]{/Users/mat/Documents/Git/LilyPond_transpose/python/tokenizer.py}

\subsection{tokens.py}
\lstinputlisting[caption={トークン定義 tokens.py}]{/Users/mat/Documents/Git/LilyPond_transpose/python/tokens.py}

\subsection{transpose.py}
\lstinputlisting[caption={移調処理 transpose.py}]{/Users/mat/Documents/Git/LilyPond_transpose/python/transpose.py}

\subsection{writer.py}
\lstinputlisting[caption={LilyPondテキスト出力 writer.py}]{/Users/mat/Documents/Git/LilyPond_transpose/python/writer.py}

\end{document}
